\documentclass{beamer}
\usepackage{beamerthemeCambridgeUS}
\usepackage{color}  % for \textcolor{red}{blah}
\usepackage{graphicx}
\usepackage{subfigure}
\usepackage{multicol}
\usepackage{verbatim} % for \begin{comment}
\usepackage{lipsum}


\newcommand\Fontvi{\fontsize{6}{7.2}\selectfont}


\newenvironment{packed_enum}{
\begin{enumerate}
  \setlength{\itemsep}{1pt}
  \setlength{\parskip}{0pt}
  \setlength{\parsep}{0pt}
}{\end{enumerate}}

\newenvironment{packed_item}{
\begin{itemize}
  \setlength{\itemsep}{1pt}
  \setlength{\parskip}{0pt}
  \setlength{\parsep}{0pt}
}{\end{itemize}}

\usetheme{CambridgeUS}

\title{Discrete Mathematics}
\author[Geoffrey Buhl]{Geoffrey Buhl \\ \texttt{geoffrey.buhl@csuci.edu}}
\institute[CSUCI]{California State University Channel Islands}
\date{week 4 class 1}

\def \vec{{\rm vec}}
\newtheorem{conjecture}[theorem]{Conjecture}
\newtheorem{explanation}[theorem]{Explanation}


\begin{document}

%\frame{\titlepage}

%\section[Outline]{}
%\frame{\tableofcontents}

\section{Week 4 Class 1}

\frame
{
  \frametitle{Introduction to Direct proofs}


Class Session Goals:
\begin{packed_enum}
\item Present Direct Proof Technique
\item Practice Direct Proofs
\end{packed_enum}
}


\frame
{
  \frametitle{What is a Direct proof?}


The simplest ,from a logic perspective, style of proof is a direct proof. Often all that is required to prove something is a systematic explanation of what everything means. Direct proofs are especially useful when proving implications. The general format to prove that statement P implies statement Q is:\vfill
Assume statement P is true. Explain, explain, …, explain. Therefore statement Q is true.
\vfill
In this course there are two sorts of flavors of direct proof, ``wordy'' and ``equation-y''.
}

\frame
{
  \frametitle{``Equation-y'' Direct Proof Example}
%\Fontvi


Here is a  ``equation-y'' or  symbolic manipulation direct proof of Pascal's Triangle Identity.
\vfill
{\bf Claim:} $\binom{n}{k}+\binom{n}{k-1}=\binom{n+1}{k}$
\pause \vfill
{\bf Proof:}
\Fontvi
\begin{eqnarray}
\binom{n}{k}+\binom{n}{k-1}
&=& \frac{n!}{k!(n-k)!}+\frac{n!}{(k-1)!((n-(k-1)!}\\
&=& \frac{n!}{k!(n-k)!}+\frac{n!}{(k-1)!((n+1-k)!}\\
&=& \frac{n+1-k}{n+1-k} \cdot \frac{n!}{k!(n-k)!}+\frac{k}{k} \cdot\frac{n!}{(k-1)!(n+1-k)!}\\
&=&  \frac{(n+1-k)n!}{k!(n+1-k)!}+\frac{k\cdot n!}{k!(n+1-k)!}\\
&=&  \frac{(n+1)n!-k\cdot n!+k\cdot n!}{k!(n+1-k)!}\\
&=& \frac{(n+1)!}{k!(n+1-k)!}\\
&=& \binom{n+1}{k} \square
\end{eqnarray}
%Thus the right-hand side of the claim is equal to the left hand-side.$\square$


}


\frame
{
  \frametitle{Second Direct Proof Example}
%\Fontvi


Here is a bit of a mixture.
\vfill
{\bf Claim:} If $n$ is an even integer then $n^2$ is even integer.
\pause \vfill
{\bf Proof:}
%\Fontvi
Let  $n$ be an even integer.   Then  $n=2k$ for some integer $k$.  Now $n^2=(2k)^2=4k^2=2(2k^2)$ where $2k^2$ is an integer.  Hence $n^2$ is even. $\square$


}

\frame
{
  \frametitle{``Wordy'' Direct Proof Example}
%\Fontvi


Here is a bit of a mixture.
\vfill
{\bf Claim:}If $n^2$ is an even integer then $n$ is an even integer.
\pause \vfill
{\bf Proof:}

Let  $n$ be an  integer.  There are two cases for $n$ it can be odd or even.  If $n$ is odd then $n^2$ is odd.  If $n$ is even then $n^2$ is even. The only case where  $n^2$ is even is the when $n$ is even. $\square$
}



\frame
{
  \frametitle{Activity}

In small groups, work on the ``Direct Proofs'' worksheet. The goals of this worksheet are to:
\vfill
\begin{packed_enum}
\item Practice creating direct proofs
\end{packed_enum}


}

\frame
{
  \frametitle{``If and only if'' Proofs}
%\Fontvi

For if any only if proofs, you need prove two parts of the proof.  If you are trying to proof the statement: $A$ if and only if $B$.  You need to prove: If $A$ then $B$ and If $B$ then $A$.
\vfill
So today we have proven:
\vfill
{\bf Claim:} $n^2$ is an even integer if and only if $n$ is an even integer.  

\vfill The two separate implications are what we have shown already:\vfill

{\bf Claim:}If $n^2$ is an even integer then $n$ is an even integer.\\

{\bf Claim:} If $n$ is an even integer then $n^2$ is even integer.

}

\begin{comment}
\frame
{
  \frametitle{Answers for Direct Proof Activity}

{\bf Problem 1}: Show $\binom{2n}{2}=2\binom{n}{2}+n^2$ \vspace{.25cm}


}


\frame
{
  \frametitle{Answers for Direct Proof Activity}


{\bf Problem 2}: Show $\binom{n}{1} ^2=\binom{n}{2}+ \binom{n+1}{2}$

}

\frame
{
  \frametitle{Answers for Direct Proof Activity}


{\bf Problem 3}: Show $\binom{n}{0}+\binom{n+1}{1}+\binom{n+2}{2} =\binom{n+3}{2}$ \vspace{.25cm}


}


\frame
{
  \frametitle{Answers for Direct Proof Activity}


{\bf Problem 4}: Show
$\binom{n}{k}\binom{n-k}{l}=\binom{n}{l}\binom{n-l}{k}$ with $k+l \geq n$.
\vspace{.25cm}


}


\frame
{
  \frametitle{Answers for Direct Proof Activity}



{\bf Problem 5}:
Let $n$ be a positive integer.  Show that $\binom{2n}{n+1}+\binom{2n}{n}=\frac{1}{2}\binom{2n+2}{n+1}$
\vspace{.25cm}



}


\frame
{
  \frametitle{Answers for Direct Proof Activity}



{\bf Problem 6}:
If $n$ an odd integer, then $n^2$ is odd.\vspace{.25cm}


}


\frame
{
  \frametitle{Answers for Direct Proof Activity}



{\bf Problem 7}:
Let $n$ any integer, then $n^3-n$ is even.\vspace{.25cm}


}


\frame
{
  \frametitle{Answers for Direct Proof Activity}


{\bf Problem 8}:
Let $a$, $b$, and $c$ be integers. If $a$ divides $b$, and $b$ divides $c$, then $a$ divides $c$.\vspace{.25cm}

}


\frame
{
  \frametitle{Answers for Direct Proof Activity}

{\bf Problem 9}:
Let $a$, $b$, and $c$ be integers. If $a^2+b^2=c^2$ is even, then at least one of the integers $a$ and $b$ is even.\vspace{.25cm}

}

\end{comment}

\end{document}

